\chapter{Installation}
\label{txt:installation}

\section{Choosing how to run Tellervo}

Tellervo is a cross-platform, open-source application for dendrochronological
measurement, data management, analysis, and curation. It can run in either full
or lite mode.

\begin{description}
 \item[Full Tellervo] connects to a Tellervo Server. The server stores series,
 metadata, users, permissions, and curatorial records in a shared PostgreSQL
 database. This is the normal choice for a laboratory or any collaborative
 installation.
 \item[Tellervo-lite] does not require a server. It reads and writes local
 dendrochronology files and provides the measuring and series-editing tools, but
 not the shared metadata, mapping, permissions, or curation facilities.
\end{description}

If your laboratory already operates a Tellervo Server, ask its administrator
for the webservice URL, your username, and your initial password. You only need
to install the desktop application. If you are responsible for a new full
Tellervo installation, install the server first and then install the desktop
application on each user's computer.

\section{Installing Tellervo Server}
\index{Installation!Server}
\label{txt:serverinstallation}

Tellervo Server 2.0 is packaged for current Debian- and Ubuntu-based systems.
The standard package installs Apache, PHP, PostgreSQL, PostGIS, PL/Java, the
Tellervo webservice, and the database administration tools. The server is not
distributed as the former VirtualBox appliance.

\subsection{Enable the PostgreSQL package repository}

PL/Java is not supplied by the main Debian and Ubuntu repositories supported
by Tellervo. Before installing Tellervo on any host that will run the database,
enable the official PostgreSQL APT repository (PGDG). The
\verb|postgresql-common| package provides a script that installs the repository
configuration and signing key:

\begin{lstlisting}[language=bash]
sudo apt update
sudo apt install postgresql-common ca-certificates
sudo /usr/share/postgresql-common/pgdg/apt.postgresql.org.sh
sudo apt update
\end{lstlisting}

The repository script is interactive and may ask for confirmation of the
detected operating-system release. It only needs to be run once on each
database host. It is not required on a separate web-only host.

Verify that APT can see the PL/Java package matching the PostgreSQL provider
you intend to install. For PostgreSQL 17:

\begin{lstlisting}[language=bash]
apt-cache policy postgresql-17-pljava
\end{lstlisting}

The output should contain a candidate version. If it shows
\verb|Candidate: (none)|, do not continue with the Tellervo database package:
check that the PGDG script completed successfully and that the operating
system can reach \verb|apt.postgresql.org|. A PostgreSQL provider cannot be
used unless its matching \verb|postgresql-MAJOR-pljava| package is available.

\subsection{Single-host installation}

Most laboratories should run the webservice and database on one server. If the
Tellervo and PostgreSQL APT repositories have been configured, install the
meta-package:

\code{sudo apt update}
\code{sudo apt install tellervo-server}

The package selects a supported PostgreSQL dependency bundle. Tellervo 2.0
packages provide PostgreSQL 18 and PostgreSQL 17 alternatives; APT selects the
first complete alternative available for the operating system.

If you received local package files instead of using an APT repository, install
the meta-package and all its local dependencies together with
\verb|apt install ./filename.deb|. Do not use \verb|dpkg --install|: it cannot
resolve the external PostgreSQL, PostGIS, PL/Java, Apache, and PHP dependencies.
For example:

\begin{lstlisting}
sudo apt install \
  ./tellervo-server-common-2.0.deb \
  ./tellervo-server-webservice-2.0.deb \
  ./tellervo-server-db-pg17-2.0.deb \
  ./tellervo-server-db-2.0.deb \
  ./tellervo-server-2.0.deb
\end{lstlisting}

After installation, run the configuration wizard and its tests:

\code{sudo tellervo-server --configure}
\code{sudo tellervo-server --test}

The wizard creates the database and database role, writes the webservice
configuration, configures Apache, applies database upgrades, and reports the
webservice URL. Save that URL for desktop users. Public DNS, HTTPS certificates,
firewalls, and reverse proxies are normally completed by the local systems
administrator after this wizard.

\subsection{Separate web and database hosts}
\label{txt:splitinstall}

Large installations may place PostgreSQL and the Apache/PHP webservice on
different hosts. The example below uses documentation-only private addresses:
the database host is \verb|192.0.2.10| and the web host is
\verb|192.0.2.20|. It creates the named instance \verb|lab-a|, with database
and login role both named \verb|tellervo_lab_a|. Replace all of these values
with the private addresses and names for your site.

\subsubsection{Prepare the database host}

First enable PGDG as described above and verify that the matching PL/Java
package is available. Then install the common package, the package matching the
PostgreSQL major version, and the database package:

\begin{lstlisting}[language=bash]
sudo apt install \
  ./tellervo-server-common-2.0.deb \
  ./tellervo-server-db-pg17-2.0.deb \
  ./tellervo-server-db-2.0.deb
\end{lstlisting}

Create the group role expected by the packaged database, a login role for this
instance, and an empty database. The \verb|\password| command prompts securely
for the password instead of placing it in shell history. The login is
currently a PostgreSQL superuser because the server's PL/Java functions and
upgrade process require elevated privileges. Protect this credential and do
not reuse it for another application.

\begin{lstlisting}[language=SQL]
sudo -u postgres psql
SET password_encryption = 'scram-sha-256';
DO $$
BEGIN
  IF NOT EXISTS (
    SELECT FROM pg_roles WHERE rolname = 'Webgroup'
  ) THEN
    CREATE ROLE "Webgroup"
      SUPERUSER INHERIT NOCREATEDB NOCREATEROLE;
  END IF;
  IF NOT EXISTS (
    SELECT FROM pg_roles WHERE rolname = 'tellervo'
  ) THEN
    CREATE ROLE tellervo NOLOGIN;
  END IF;
  IF NOT EXISTS (
    SELECT FROM pg_roles WHERE rolname = 'pbrewer'
  ) THEN
    CREATE ROLE pbrewer NOLOGIN;
  END IF;
END
$$;
CREATE ROLE tellervo_lab_a
  LOGIN SUPERUSER INHERIT NOCREATEDB NOCREATEROLE;
GRANT "Webgroup" TO tellervo_lab_a;
CREATE DATABASE tellervo_lab_a;
\password tellervo_lab_a
\quit
\end{lstlisting}

Use fresh database and login-role names. If either already exists, stop and
verify that it is not an installation you need to preserve. The non-login
\verb|tellervo| and \verb|pbrewer| roles are retained only because some
packaged upgrade scripts refer to these historical object owners; they cannot
be used to connect.

Initialise the database with the PL/Java setup and packaged database template:

\begin{lstlisting}[language=bash]
sudo -u postgres psql \
  --dbname=tellervo_lab_a \
  --set=ON_ERROR_STOP=1 \
  --file=/usr/share/tellervo-server/db-upgrade-patches/database_upgrade-1.3.0e.notransaction.sql

sudo -u postgres pg_restore \
  --exit-on-error \
  --no-owner \
  --dbname=tellervo_lab_a \
  /usr/share/tellervo-server/db-templates/tellervo_database_template_2.0.sql
\end{lstlisting}

Despite its \verb|.sql| suffix, the template is a PostgreSQL custom-format
archive and must be loaded with \verb|pg_restore|.

\subsubsection{Permit the web host to connect}

Ask PostgreSQL for its active configuration paths rather than assuming a
version-specific directory:

\begin{lstlisting}[language=bash]
sudo -u postgres psql -Atqc "SHOW config_file;"
sudo -u postgres psql -Atqc "SHOW hba_file;"
\end{lstlisting}

Edit the reported \verb|postgresql.conf| and bind PostgreSQL to the database
host's private address as well as localhost:

\begin{lstlisting}
listen_addresses = 'localhost,192.0.2.10'
\end{lstlisting}

Edit the reported \verb|pg_hba.conf| and add this line before any broader
reject rule:

\begin{lstlisting}
host  tellervo_lab_a  tellervo_lab_a  192.0.2.20/32  scram-sha-256
\end{lstlisting}

This permits only the named database and role, from the single web-host
address, using SCRAM password authentication. Do not use an unrestricted rule
such as \verb|host all all 0.0.0.0/0|, and do not expose PostgreSQL to the
public Internet. Restart PostgreSQL because a \verb|listen_addresses| change
cannot be applied by a reload alone:

\begin{lstlisting}[language=bash]
sudo systemctl restart postgresql
\end{lstlisting}

If a host or network firewall is enabled, permit TCP port 5432 only from the
web host. For example, with UFW:

\begin{lstlisting}[language=bash]
sudo ufw allow from 192.0.2.20 \
  to 192.0.2.10 port 5432 proto tcp
\end{lstlisting}

\subsubsection{Configure and test the web host}

Install the web-service packages and PostgreSQL client on the web host:

\begin{lstlisting}[language=bash]
sudo apt install postgresql-client \
  ./tellervo-server-common-2.0.deb \
  ./tellervo-server-webservice-2.0.deb
\end{lstlisting}

First test the network route, authentication, and HBA rule independently of
Tellervo. Enter the password created above when prompted:

\begin{lstlisting}[language=bash]
psql --host=192.0.2.10 --port=5432 \
  --dbname=tellervo_lab_a \
  --username=tellervo_lab_a \
  --command='SELECT current_database(), current_user;'
\end{lstlisting}

Once this succeeds, run the webservice wizard and supply the same database
host, port, database name, role, and password:

\code{sudo tellervo-server --instance lab-a --configure}
\code{sudo tellervo-server --instance lab-a --test}

The wizard deliberately does not create a database on a remote PostgreSQL
host. It verifies the prepared database and applies any outstanding upgrades.
See chapter \ref{txt:runningserver} for named instances and routine
administration.

\section{Installing the desktop application}
\label{txt:desktopinstall}
\index{Installation!Desktop application}

Tellervo Desktop 2.0 uses Java 21 and is distributed as a native package
containing its Java 21 runtime. Users of the native packages do not need to
install Java separately. Obtain the package for your operating system from the
Tellervo download page or the project's release assets. A Java 21 JDK is
required only when building Tellervo or running it directly from Maven or a
JAR.

\begin{description}
 \item[Windows] Run the MSI package and follow Windows Installer. The default
 per-user installation is placed below \verb|%LOCALAPPDATA%\Tellervo| and does
 not require administrator privileges. The required 64-bit serial library is
 included.
 \item[macOS] Open the DMG and install Tellervo in the usual Applications
 location. Depending on local security policy, the first launch may require
 approval in macOS Privacy \& Security settings.
 \item[Debian or Ubuntu Linux] Install the DEB with APT so that the RXTX and
 JOGL runtime dependencies are installed:

 \code{sudo apt install ./tellervo\_2.0.0\_amd64.deb}

 The package adds an application-menu entry and a \verb|tellervo| command in
 \verb|/usr/bin|.
\end{description}

Current Tellervo desktop packages are 64-bit. Measuring hardware connected by
USB or a serial adapter may require an operating-system driver and permission
for the user to access the device. See chapter
\ref{txt:MeasuringPlatformConfig}.

\subsection{First launch}
\index{Wizard, Setup}

On first launch, Tellervo opens its setup wizard. The wizard can be run again
from \menutwo{Help}{Setup Wizard}; the same settings can be changed in
\menutwo{Edit}{Preferences}.

\begin{description}
 \item[Tellervo mode] Choose full mode to connect to a server or lite mode to
 work with local files.
 \item[Network connection] In most environments the system proxy settings are
 correct. Enter a manual proxy only when instructed by your network
 administrator.
 \item[Tellervo Server] In full mode, enter the complete webservice URL supplied
 by the server administrator. Use HTTPS for servers reached across an
 untrusted network.
 \item[Measuring platform] Select the connected device and port, enter any
 device-specific serial settings, and use Test Connection before measuring
 samples.
\end{description}

Full mode then displays the Tellervo login dialog. These are Tellervo database
credentials, not operating-system credentials. Saving a password is appropriate
only on a trusted personal workstation.

\subsection{Mapping support}
\index{Mapping}

Tellervo's three-dimensional map requires working OpenGL support. Native
packages include the Java graphics components, but the operating system must
provide a suitable graphics driver. Open \menutwo{Admin}{Site map} to test map
initialization. If it fails, update the graphics driver and consult the system
information and error log under the Help menu.

\section{Upgrading}
\index{Upgrading}
\index{Upgrading!Tellervo desktop}

To upgrade Tellervo Desktop, close the application and install the new native
package over the existing version. User preferences are retained. If a release
requires a newer Tellervo Server, upgrade and test the server before deploying
the desktop package to users.

Always back up the server before an upgrade. A server installed from an APT
repository is upgraded with:

\code{sudo apt update}
\code{sudo apt upgrade}
\code{sudo tellervo-server --test}

Local DEB upgrades should again use \verb|apt install| with all required local
packages. Detailed backup, database upgrade, and multi-instance procedures are
in chapter \ref{txt:runningserver}.

\section{Uninstalling}

Use Windows Installed Apps, remove the macOS application, or use the Linux
package manager to uninstall Tellervo Desktop. On Debian or Ubuntu:

\code{sudo apt remove tellervo}

\warn{Removing server packages is not the same as deleting a Tellervo
database. Before changing or decommissioning a server, make a database backup
and a separate backup of uploaded media files. Verify both backups on another
machine.}

To remove the server software while preserving data, use:

\code{sudo apt remove tellervo-server}

Do not purge packages, delete a named instance, or remove PostgreSQL data
directories unless the database and media are intentionally being destroyed.
The named-instance deletion command is interactive because it removes both the
instance database and its web folder.
